\documentclass[a4paper,11pt]{article}
\usepackage[T1]{fontenc}
\usepackage[utf8]{inputenc}
\usepackage[ngerman]{babel}
\usepackage[margin=2.5cm]{geometry}
\usepackage{booktabs}

\begin{document}

\section*{Blatt 3, Aufgaben 3.2.3 und 3.2.4}

Für jede Kombination aus Verteilung, Variationskoeffizient (CV) und
Stichprobengröße wurden 500 Stichproben mit wahrem Mittelwert \(\mu=10\)
simuliert. Für jede Stichprobe berechnet \texttt{study.DCCTest} beide
Konfidenzintervalle aus denselben Ziehungen. Die Tabellen zeigen, wie oft
die Intervalle den Wert 10 enthielten. Pro Zelle steht zuerst die
Überdeckung des 90\,\%-Intervalls, dann die des 95\,\%-Intervalls, jeweils
in Prozent. Der Zufallszahlengenerator verwendet den festen Startwert 42.

\subsection*{Normalverteilung}

\begin{center}
\begin{tabular}{@{}rcccc@{}}
\toprule
CV & \(n=5\) & \(n=10\) & \(n=50\) & \(n=100\) \\
\midrule
0,25 & 91,0 / 95,4 & 89,6 / 94,0 & 91,8 / 96,0 & 90,6 / 94,6 \\
0,5  & 90,8 / 95,0 & 88,8 / 93,0 & 90,0 / 93,6 & 92,0 / 95,8 \\
1    & 90,8 / 94,4 & 89,2 / 95,0 & 91,8 / 95,0 & 91,4 / 96,2 \\
2    & 91,2 / 96,4 & 89,4 / 94,6 & 91,4 / 96,4 & 91,8 / 97,2 \\
4    & 90,0 / 95,4 & 89,2 / 96,4 & 88,4 / 95,2 & 89,8 / 95,6 \\
\bottomrule
\end{tabular}
\end{center}

Schon bei \(n=5\) liegen die Trefferquoten nahe bei den Sollwerten. Das
passt zum \(t\)-Intervall für normalverteilte Daten. Ein höherer CV macht
die Intervalle breiter. Die theoretische Trefferquote bleibt dabei gleich,
weil die Streuung sowohl die Abweichung des Stichprobenmittels als auch
den geschätzten Standardfehler im selben Verhältnis vergrößert. Bei 500
Wiederholungen beträgt der Monte-Carlo-Standardfehler nahe den Sollwerten
rund 1,34 bzw. 0,97 Prozentpunkte.

\subsection*{Erlang-, Exponential- und Hyperexponentialverteilung}

\begin{center}
\begin{tabular}{@{}lrcccc@{}}
\toprule
Verteilung & CV & \(n=5\) & \(n=10\) & \(n=50\) & \(n=100\) \\
\midrule
Erlang (\(k=16\)) & 0,25 & 89,6 / 94,6 & 90,8 / 94,4 & 89,4 / 95,6 & 88,6 / 95,2 \\
Erlang (\(k=4\))  & 0,5  & 85,2 / 91,8 & 89,6 / 94,4 & 90,0 / 93,8 & 90,0 / 95,4 \\
Erlang (\(k=1\))  & 1    & 83,0 / 88,0 & 86,0 / 90,6 & 89,6 / 93,6 & 88,2 / 94,0 \\
Exponential      & 1    & 80,8 / 87,2 & 84,2 / 88,8 & 86,6 / 92,6 & 87,8 / 94,0 \\
\(H_2\), gleiche Raten & 1 & 83,8 / 91,0 & 83,0 / 87,8 & 87,6 / 92,2 & 91,0 / 94,8 \\
\(H_2\)          & 2    & 64,8 / 71,8 & 65,2 / 72,4 & 82,4 / 86,0 & 87,0 / 91,0 \\
\(H_2\)          & 4    & 46,6 / 52,8 & 39,6 / 45,2 & 62,2 / 64,2 & 72,8 / 76,2 \\
\bottomrule
\end{tabular}
\end{center}

Bei Erlang gilt \(\mathrm{CV}=1/\sqrt{k}\), daher passen aus der Vorgabe
nur \(k=16,4,1\). Die Exponentialverteilung hat immer CV~1.
Für \(H_2\) gilt \(\mathrm{CV}\geq1\). Bei CV~1 ergeben gleiche Raten
ebenfalls eine Exponentialverteilung. Die drei Zeilen mit dieser
Verteilung beruhen auf verschiedenen Ziehungen und fallen deshalb
etwas unterschiedlich aus.

Die Rechtsschiefe macht sich vor allem bei kleinen Stichproben bemerkbar.
Bei der Exponentialverteilung enthält das 95\,\%-Intervall für \(n=5\)
den Mittelwert nur in 87,2\,\% der Fälle. Bei \(n=100\) sind es 94,0\,\%.
Für \(H_2\) mit CV~4 bleibt der Abstand deutlich größer. Die langsame Phase
tritt pro Ziehung nur mit rund 3\,\% Wahrscheinlichkeit auf und erzeugt
besonders große Werte. Kleine Stichproben erfassen diese Phase oft gar
nicht. Dann können sowohl der geschätzte Mittelwert als auch der
geschätzte Standardfehler zu klein ausfallen. Selbst bei \(n=100\)
enthalten hier nur 76,2\,\% der 95\,\%-Intervalle den Mittelwert.

\end{document}
